%% ===========================================================================
%%  Bab 2 Dasar Pemikiran
%%  Aturan isi: alinea 4-5 kalimat, satu gagasan per alinea, spasi 2, awal
%%  alinea menjorok otomatis 1,5 cm, rata kiri-kanan. Istilah asing italic:
%%  \asing{retrieval-augmented generation}.
%%  Sudut pandang: orang pertama ("saya") untuk pengalaman magang.
%% ===========================================================================
\chapter{Dasar Pemikiran}

\section{Dasar Pemikiran Pemilihan Lokasi Magang}

Keinginan memperoleh pengalaman belajar di luar negeri telah muncul sebelum
saya memasuki tahap akhir studi. Saya berkeinginan mengikuti kegiatan akademik
di negara lain untuk mengamati cara belajar dan cara kerja yang berbeda.
Harapan itu tumbuh karena bidang sistem informasi dipakai pada berbagai jenis
organisasi dengan kebiasaan kerja yang tidak seragam. Pengalaman lintas negara
dinilai dapat memperluas wawasan mengenai penerapan bidang studi pada konteks
yang berbeda.

Keinginan tersebut belum memperoleh jalan karena informasi mengenai program
lintas negara masih terbatas. Saya pada awalnya memahami bahwa pengalaman
belajar di luar negeri hanya tersedia melalui jalur pertukaran mahasiswa yang
membutuhkan waktu dan biaya tersendiri. Manajemen \NamaKampus{} kemudian
menyampaikan informasi mengenai program magang studi luar negeri yang dapat
diikuti oleh mahasiswa. Informasi tersebut menjelaskan ketentuan pendaftaran
dan kewajiban peserta selama mengikuti program. Informasi itu membuka pilihan
yang lebih terarah karena kegiatan magang dapat diselesaikan sekaligus
memenuhi kewajiban studi.

Setelah mempelajari ketentuan program, saya memutuskan untuk mendaftar. Saya
menyiapkan berkas yang diminta dan mengikuti proses pendaftaran sampai
memperoleh keputusan penerimaan. Saya dinyatakan diterima di Broadband
Multimedia Wireless Laboratory, National Taiwan University of Science and
Technology. Penerimaan tersebut menjadi titik awal pelaksanaan magang pada
periode yang telah ditetapkan. Saya kemudian mempelajari gambaran pekerjaan
laboratorium agar dapat menyesuaikan kemampuan dengan kebutuhan yang akan
dihadapi.

Penyesuaian kemampuan diperlukan karena Broadband Multimedia Wireless
Laboratory memiliki kegiatan yang lebih luas daripada pengembangan jaringan
nirkabel. Laboratorium juga mengembangkan perangkat lunak pendukung penelitian,
seperti pengelolaan infrastruktur komputasi, pengolahan rekaman rapat,
inventaris peralatan, dan alat bantu penulisan artikel. Ragam pekerjaan
tersebut membutuhkan analisis kebutuhan, perancangan data, dan pengembangan
antarmuka. Ketiga kemampuan itu merupakan bagian dari kompetensi Program Studi
Sistem Informasi. Dengan demikian, penempatan di laboratorium memberi saya
kesempatan menerapkan bidang studi pada pekerjaan yang digunakan peneliti.

Lingkungan laboratorium juga memberikan pengalaman bekerja bersama anggota dari
latar belakang yang berbeda. Diskusi dilakukan dalam bahasa Inggris dan
mengikuti kebiasaan kerja tim peneliti, termasuk rapat mingguan serta
pemeriksaan kemajuan secara berkala. Saya memperoleh kesempatan menjelaskan
rancangan, menerima masukan, dan memperbaiki pekerjaan berdasarkan masukan
tersebut. Kebiasaan itu melatih cara menyampaikan gagasan secara ringkas dan
menanggapi pertanyaan secara terbuka. Pengalaman tersebut menjadi bekal untuk
bekerja pada proyek yang melibatkan peserta dari beberapa negara.

Pertimbangan terakhir adalah ketersediaan bukti pekerjaan untuk penyusunan
laporan magang. Setiap proyek di laboratorium memiliki catatan kegiatan,
repositori, dan dokumentasi yang dapat ditelusuri kembali. Saya dapat
menghubungkan keputusan teknis dengan kebutuhan pengguna yang
melatarbelakanginya. Hubungan itu memudahkan penyusunan pembahasan tanpa
mengandalkan ingatan atau keterangan lisan. Dengan demikian, lokasi magang
menyediakan bahan yang cukup untuk dinilai secara akademik.

\section{Dasar Pemikiran Pemilihan Bidang Magang}

Pada awal magang, saya ditempatkan pada pekerjaan infrastruktur komputasi. Saya
menyiapkan klaster, otomasi penyebaran, dan pemantauan layanan yang menopang
percobaan laboratorium. Pekerjaan itu saya jalani sebagai bidang utama karena
sesuai dengan keterampilan teknis yang saya pelajari di kampus. Setelah saya
dinyatakan diterima dan mengikuti pekerjaan laboratorium lebih lama, saya
mengetahui adanya persoalan lain yang ternyata lebih mendesak. Persoalan itu
muncul pada kegiatan penulisan artikel jurnal yang menjadi pekerjaan utama
laboratorium.

Laboratorium membutuhkan cara agar penulisan artikel berjalan lebih efektif,
terutama pada tahap pencarian rujukan. Peneliti harus menemukan artikel yang
sesuai dengan topik, lalu menilai apakah metode dan temuannya mendukung
penelitian yang sedang dikerjakan. Pencarian dengan kata kunci menghasilkan
banyak artikel yang belum tentu sesuai dengan konteks penelitian. Akibatnya,
peneliti harus membaca dan menyaring abstrak satu per satu. Kebutuhan itulah
yang kemudian mengalihkan bidang magang saya dari infrastruktur ke pencarian
data riset.

Kecerdasan buatan menawarkan kemungkinan membantu peneliti membaca dan menilai
bahan dalam jumlah besar. \citet{munawar2023manfaat} menunjukkan bahwa model
bahasa dapat membantu peneliti menyusun bahan dan draf awal, dengan catatan
hasilnya tetap diperiksa. Namun, \citet{ariqo2025analisis} mengingatkan bahwa
penggunaan kecerdasan buatan dalam publikasi ilmiah membawa tantangan etika dan
ketepatan. Model bahasa dapat menghasilkan pernyataan yang terdengar meyakinkan
tanpa sumber yang benar. Karena itu, bidang ini membutuhkan rancangan yang
menempatkan pemeriksaan manusia sebagai bagian dari sistem.

Pendekatan yang saya pilih adalah pencocokan konteks penelitian dengan abstrak
artikel melalui \asing{retrieval-augmented generation}. Pendekatan ini
mengambil dokumen lebih dahulu, lalu menyusun jawaban hanya berdasarkan dokumen
yang diambil. \citet{lewis2020retrieval} merancang cara kerja tersebut, sedangkan
\citet{budi2024repositori} menerapkannya pada repositori institusi dan
\citet{tohir2024utilizing} mengujinya pada korpus berbahasa Indonesia. Kedua
penerapan itu menunjukkan bahwa sumber yang diambil lebih dahulu membuat
jawaban dapat ditelusuri. Dengan cara itu, artikel yang ditampilkan lebih dekat
dengan konteks penelitian, bukan sekadar cocok pada kata kunci.

Bidang ini juga sesuai dengan kompetensi Program Studi Sistem Informasi.
Pekerjaan tidak berhenti pada pemilihan model bahasa, tetapi mencakup analisis
kebutuhan, perancangan basis data, pengembangan antarmuka, integrasi layanan,
dan pengujian. Saya harus memahami data artikel, menyusun alur pencarian, serta
menentukan informasi yang ditampilkan kepada pengguna. Keterampilan tersebut
merupakan bagian dari sistem informasi dan dapat diuji pada pemakaian nyata.
Dengan demikian, peralihan bidang magang tetap berada dalam lingkup studi saya.

Peralihan tersebut juga memberi manfaat jangka panjang bagi saya. Kemampuan
merancang sistem yang menggunakan model bahasa secara terkendali semakin
dibutuhkan pada pekerjaan pengembangan perangkat lunak. Saya tidak hanya belajar
memanggil layanan kecerdasan buatan, tetapi juga belajar menilai batasannya.
Pengalaman itu menjadi bekal untuk bekerja pada sistem informasi yang
membutuhkan ketepatan sumber dan tanggung jawab pengguna.
